\documentclass[11pt,a4paper]{article}

\usepackage[T1]{fontenc}
\usepackage[utf8]{inputenc}
\usepackage[french]{babel}
\usepackage{amsmath,amssymb,amsthm}
\usepackage{mathtools}
\usepackage[margin=2.6cm]{geometry}
\usepackage{enumitem}
\usepackage{xcolor}
\usepackage[colorlinks=true, linkcolor=blue!60!black, citecolor=blue!60!black]{hyperref}

% ----- Environnements de théorèmes -----
\theoremstyle{plain}
\newtheorem{theorem}{Théorème}[section]
\newtheorem{proposition}[theorem]{Proposition}
\newtheorem{lemma}[theorem]{Lemme}
\newtheorem{corollary}[theorem]{Corollaire}
\theoremstyle{definition}
\newtheorem{definition}[theorem]{Définition}
\newtheorem{example}[theorem]{Exemple}
\theoremstyle{remark}
\newtheorem{remark}[theorem]{Remarque}

% ----- Macros -----
\newcommand{\E}{\mathbb{E}}
\newcommand{\Prob}{\mathbb{P}}
\newcommand{\R}{\mathbb{R}}
\newcommand{\Q}{\mathbb{Q}}
\newcommand{\N}{\mathbb{N}}
\newcommand{\F}{\mathcal{F}}
\newcommand{\G}{\mathcal{G}}
\newcommand{\ps}{\text{p.s.}}
\newcommand{\ind}{\mathbf{1}}
\DeclarePairedDelimiter{\abs}{\lvert}{\rvert}
\DeclarePairedDelimiter{\norm}{\lVert}{\rVert}

\title{Convergence des martingales de Doob,\\
uniforme intégrabilité\\
et dichotomie de Kakutani}
\author{}
\date{Juillet 2026}

\begin{document}
\maketitle

\begin{abstract}
Ce document présente le théorème de convergence presque sûre des martingales de
Doob avec la preuve complète de l'inégalité des montées, développe en détail la
notion d'uniforme intégrabilité (critères, théorème de Vitali, critère de
La~Vallée Poussin), caractérise les martingales fermées et le théorème d'arrêt
associé, puis établit la dichotomie de Kakutani (1948) pour les martingales
produit et les mesures produit infinies. Le cas gaussien est traité complètement,
avec le théorème de Cameron--Martin en corollaire, et les liens avec la
statistique asymptotique (contiguïté de Le~Cam) et la finance mathématique
(Girsanov, Novikov) sont détaillés.
\end{abstract}

\tableofcontents

%=======================================================================
\section{Le théorème de convergence de Doob}
%=======================================================================

Dans toute la suite, $(\Omega, \F, (\F_n)_{n \ge 0}, \Prob)$ désigne un espace de
probabilité filtré. Un processus $(X_n)_{n\ge0}$ adapté et intégrable est une
\emph{sous-martingale} si $\E[X_{n+1} \mid \F_n] \ge X_n$ p.s.\ pour tout $n$,
une \emph{surmartingale} si l'inégalité inverse a lieu, une \emph{martingale} si
l'égalité a lieu.

\begin{theorem}[Doob]\label{thm:doob}
Soit $(X_n)_{n\ge 0}$ une sous-martingale telle que
\begin{equation}\label{eq:hyp}
\sup_{n\ge 0} \E\bigl[X_n^+\bigr] < \infty.
\end{equation}
Alors il existe une variable aléatoire $X_\infty$ telle que
\[
X_n \xrightarrow[n\to\infty]{\ps} X_\infty,
\qquad \E\bigl[\abs{X_\infty}\bigr] < \infty.
\]
\end{theorem}

\begin{remark}\label{rem:L1-equiv}
Pour une sous-martingale, la condition \eqref{eq:hyp} équivaut à la bornitude
dans $L^1$. En effet, $\abs{x} = 2x^+ - x$ donne
$\E[\abs{X_n}] = 2\,\E[X_n^+] - \E[X_n]$, et la suite $(\E[X_n])_n$ est
croissante (prendre l'espérance dans la propriété de sous-martingale), donc
minorée par $\E[X_0]$ : ainsi
$\E[\abs{X_n}] \le 2\,\E[X_n^+] - \E[X_0]$. La réciproque est immédiate puisque
$X_n^+ \le \abs{X_n}$.
\end{remark}

\begin{corollary}\label{cor:cas-part}
Les processus suivants convergent presque sûrement vers une limite intégrable :
\begin{enumerate}[label=(\roman*)]
\item toute martingale bornée dans $L^1$ ;
\item toute sous-martingale majorée par une constante ;
\item toute surmartingale positive --- sans aucune hypothèse d'intégrabilité
      supplémentaire. En effet, $(-X_n)$ est une sous-martingale négative, donc
      $(-X_n)^+ = 0$ et \eqref{eq:hyp} est trivialement vérifiée ; de plus
      $\E[\abs{X_n}] = \E[X_n] \le \E[X_0]$.
\end{enumerate}
\end{corollary}

%-----------------------------------------------------------------------
\subsection{Transformées de martingales}
%-----------------------------------------------------------------------

L'outil de la preuve est la \emph{transformée de martingale}, version discrète de
l'intégrale stochastique. Un processus $(C_n)_{n\ge1}$ est \emph{prévisible} si
$C_n$ est $\F_{n-1}$-mesurable pour tout $n \ge 1$.

\begin{definition}
Pour $(C_n)$ prévisible borné et $(X_n)$ adapté intégrable, la transformée
$C \bullet X$ est le processus
\[
(C \bullet X)_n = \sum_{k=1}^{n} C_k\,(X_k - X_{k-1}),
\qquad (C \bullet X)_0 = 0.
\]
\end{definition}

L'interprétation de jeu est utile : $X_k - X_{k-1}$ est le gain unitaire de la
$k$-ième partie, $C_k$ la mise, décidée \emph{avant} de connaître l'issue
(prévisibilité).

\begin{lemma}\label{lem:transform}
Si $(X_n)$ est une sous-martingale et $(C_n)$ est prévisible avec
$0 \le C_n \le K$, alors $(C \bullet X)$ est une sous-martingale ; en
particulier $\E[(C\bullet X)_n] \ge 0$. Si $(X_n)$ est une martingale et
$(C_n)$ est prévisible borné (de signe quelconque), alors $(C \bullet X)$ est
une martingale.
\end{lemma}

\begin{proof}
Par prévisibilité de $C_{n+1}$,
\[
\E\bigl[(C\bullet X)_{n+1} - (C\bullet X)_n \,\big|\, \F_n\bigr]
= C_{n+1}\, \E\bigl[X_{n+1} - X_n \mid \F_n\bigr] \ge 0,
\]
le produit étant licite car $C_{n+1}$ est $\F_n$-mesurable et borné. Le signe
vient de $C_{n+1} \ge 0$ et de la propriété de sous-martingale. Le cas
martingale est identique avec une égalité.
\end{proof}

Moralité : \emph{on ne peut pas battre un jeu défavorable ou équilibré par une
stratégie de mise prévisible}. La preuve de l'inégalité des montées consiste à
exhiber une stratégie qui gagnerait si le processus oscillait trop.

%-----------------------------------------------------------------------
\subsection{L'inégalité des montées}
%-----------------------------------------------------------------------

Pour $a < b$ réels, on définit récursivement les temps d'arrêt
\[
S_1 = \inf\{n \ge 0 : X_n \le a\}, \qquad
T_j = \inf\{n \ge S_j : X_n \ge b\}, \qquad
S_{j+1} = \inf\{n \ge T_j : X_n \le a\},
\]
avec la convention $\inf \emptyset = +\infty$, et le \emph{nombre de montées}
avant l'instant $n$ :
\[
U_n[a,b] = \max\{j \ge 0 : T_j \le n\}.
\]

\begin{lemma}[Inégalité des montées de Doob]\label{lem:upcrossing}
Si $(X_n)$ est une sous-martingale, alors pour tous $a<b$ et tout $n \ge 0$,
\[
(b-a)\, \E\bigl[U_n[a,b]\bigr]
\;\le\;
\E\bigl[(X_n - a)^+\bigr] - \E\bigl[(X_0 - a)^+\bigr].
\]
\end{lemma}

\begin{proof}
\emph{Étape 1 : réduction à un processus positif.} Posons
$Y_n = (X_n - a)^+$. La fonction $x \mapsto (x-a)^+$ étant convexe croissante,
l'inégalité de Jensen conditionnelle montre que $(Y_n)$ est encore une
sous-martingale :
\[
\E[Y_{n+1} \mid \F_n] = \E\bigl[(X_{n+1}-a)^+ \mid \F_n\bigr]
\ge \bigl(\E[X_{n+1} \mid \F_n] - a\bigr)^+ \ge (X_n - a)^+ = Y_n.
\]
De plus, les montées de $X$ sur $[a,b]$ sont exactement les montées de $Y$ sur
$[0,\, b-a]$ : $X_n \le a \iff Y_n = 0$ et $X_n \ge b \iff Y_n \ge b-a$.

\emph{Étape 2 : la stratégie \og acheter bas, vendre haut \fg.} Définissons le
processus de mises $C_k \in \{0,1\}$ par $C_1 = \ind_{\{Y_0 = 0\}}$ et
\[
C_k = \ind_{\{C_{k-1} = 1\}} \ind_{\{Y_{k-1} < b-a\}}
    + \ind_{\{C_{k-1} = 0\}} \ind_{\{Y_{k-1} = 0\}},
\qquad k \ge 2 :
\]
on mise à partir du moment où $Y$ touche $0$, on cesse dès que $Y$ atteint
$b-a$, et on attend le retour en $0$ pour recommencer. Le processus $(C_k)$ est
prévisible ($C_k$ ne dépend que de $Y_0, \dots, Y_{k-1}$).

\emph{Étape 3 : minoration trajectorielle du gain.} Chaque montée complète de
$Y$ de $0$ à $b-a$ rapporte au moins $b-a$ à la stratégie. Après la dernière
montée complète, deux cas : soit la stratégie est inactive ($C = 0$), gain nul ;
soit elle est active sur une montée inachevée, et le gain partiel est
$Y_n - Y_{\tau} = Y_n \ge 0$, où $\tau$ est le dernier instant de passage en $0$
--- ce gain résiduel est positif \emph{parce que $Y \ge 0$ et que l'on achète en
$Y_\tau = 0$}. D'où, trajectoire par trajectoire,
\[
(C \bullet Y)_n \;\ge\; (b-a)\, U_n[a,b].
\]

\emph{Étape 4 : majoration en espérance.} Le processus $D_k = 1 - C_k$ est
prévisible à valeurs dans $[0,1]$, donc par le lemme~\ref{lem:transform},
$\E[(D \bullet Y)_n] \ge 0$. Or $(C\bullet Y)_n + (D \bullet Y)_n = Y_n - Y_0$,
d'où
\[
(b-a)\,\E\bigl[U_n[a,b]\bigr] \le \E\bigl[(C\bullet Y)_n\bigr]
= \E[Y_n - Y_0] - \E\bigl[(D\bullet Y)_n\bigr]
\le \E[Y_n] - \E[Y_0],
\]
qui est l'inégalité annoncée.
\end{proof}

%-----------------------------------------------------------------------
\subsection{Preuve du théorème de convergence}
%-----------------------------------------------------------------------

\begin{proof}[Preuve du théorème~\ref{thm:doob}]
Fixons $a < b$ rationnels. Comme $(X_n - a)^+ \le X_n^+ + \abs{a}$, l'hypothèse
\eqref{eq:hyp} et le lemme~\ref{lem:upcrossing} donnent, pour tout $n$,
\[
\E\bigl[U_n[a,b]\bigr] \le \frac{\sup_m \E[X_m^+] + \abs{a}}{b-a} < \infty.
\]
La suite $U_n[a,b] \uparrow U_\infty[a,b]$ étant croissante, la convergence
monotone donne $\E[U_\infty[a,b]] < \infty$, donc
$U_\infty[a,b] < \infty$ presque sûrement.

Or, si une suite réelle $(x_n)$ vérifie $\liminf x_n < \limsup x_n$, il existe
des rationnels $a < b$ intercalés, et la suite traverse $[a,b]$ en montant une
infinité de fois. Ainsi
\[
\Bigl\{\liminf_n X_n < \limsup_n X_n\Bigr\}
\subseteq \bigcup_{\substack{a<b \\ a,b \in \Q}} \bigl\{U_\infty[a,b] = \infty\bigr\},
\]
union \emph{dénombrable} d'événements négligeables, donc négligeable. La limite
$X_\infty = \lim_n X_n$ existe donc p.s.\ dans $\overline{\R} = [-\infty,
+\infty]$.

Il reste à exclure les valeurs infinies. Par le lemme de Fatou et la
remarque~\ref{rem:L1-equiv},
\[
\E\bigl[\abs{X_\infty}\bigr]
= \E\Bigl[\liminf_n \abs{X_n}\Bigr]
\le \liminf_n \E\bigl[\abs{X_n}\bigr]
\le \sup_n \E\bigl[\abs{X_n}\bigr] < \infty,
\]
donc $\abs{X_\infty} < \infty$ p.s.\ et $X_\infty \in L^1$.
\end{proof}

\begin{remark}
La preuve montre au passage que l'hypothèse \eqref{eq:hyp} ne sert qu'à borner
uniformément le membre de droite de l'inégalité des montées : c'est la
non-oscillation qui est au c\oe ur du théorème. La convergence obtenue est
presque sûre ; savoir si elle a lieu dans $L^1$ est une question strictement
plus fine, qui fait l'objet des deux sections suivantes.
\end{remark}

%=======================================================================
\section{Uniforme intégrabilité}
%=======================================================================

%-----------------------------------------------------------------------
\subsection{Définition et premières propriétés}
%-----------------------------------------------------------------------

\begin{definition}\label{def:UI}
Une famille $(X_i)_{i \in I}$ de variables aléatoires est \emph{uniformément
intégrable} (UI) si
\[
\lim_{c \to \infty} \; \sup_{i \in I} \;
\E\bigl[\abs{X_i}\, \ind_{\{\abs{X_i} > c\}}\bigr] = 0.
\]
\end{definition}

Pour une variable $X$ intégrable \emph{isolée}, on a toujours
$\E[\abs{X}\ind_{\{\abs{X}>c\}}] \to 0$ par convergence dominée
(domination par $\abs{X}$). La définition~\ref{def:UI} exige que ce contrôle des
queues soit \emph{uniforme sur la famille} : un même seuil $c$ doit fonctionner
pour tous les $X_i$ simultanément. Intuitivement : \emph{aucune masse $L^1$ ne
peut s'échapper vers l'infini le long de la famille}.

\begin{proposition}\label{prop:UI-L1}
Si $(X_i)_{i\in I}$ est UI, alors $\sup_i \E[\abs{X_i}] < \infty$. La réciproque
est fausse.
\end{proposition}

\begin{proof}
Choisissons $c$ tel que $\sup_i \E[\abs{X_i}\ind_{\{\abs{X_i}>c\}}] \le 1$.
Alors pour tout $i$,
\[
\E[\abs{X_i}] = \E\bigl[\abs{X_i}\ind_{\{\abs{X_i}\le c\}}\bigr]
+ \E\bigl[\abs{X_i}\ind_{\{\abs{X_i}>c\}}\bigr] \le c + 1.
\]
Contre-exemple pour la réciproque : $X_n = n\,\ind_{A_n}$ avec
$\Prob(A_n) = 1/n$. Alors $\E[\abs{X_n}] = 1$ pour tout $n$, mais pour tout $c$
fixé et tout $n > c$, $\E[\abs{X_n}\ind_{\{\abs{X_n}>c\}}] = 1$ : la masse
unité est intégralement portée par des valeurs qui divergent. C'est le schéma
exact de la martingale exponentielle de l'exemple~\ref{ex:gauss} ci-dessous.
\end{proof}

Une reformulation équivalente, souvent plus maniable, sépare bornitude et
\emph{absolue continuité uniforme}.

\begin{proposition}\label{prop:UI-eps-delta}
$(X_i)_{i\in I}$ est UI si et seulement si les deux conditions suivantes sont
réunies :
\begin{enumerate}[label=(\alph*)]
\item $\sup_i \E[\abs{X_i}] < \infty$ ;
\item pour tout $\varepsilon > 0$, il existe $\delta > 0$ tel que pour tout
      événement $A$ avec $\Prob(A) \le \delta$, on ait
      $\sup_i \E[\abs{X_i}\ind_A] \le \varepsilon$.
\end{enumerate}
\end{proposition}

\begin{proof}
\emph{UI $\Rightarrow$ (a), (b).} (a) est la proposition~\ref{prop:UI-L1}. Pour
(b) : soit $\varepsilon > 0$ et $c$ tel que
$\sup_i \E[\abs{X_i}\ind_{\{\abs{X_i}>c\}}] \le \varepsilon/2$ ; alors pour
$\Prob(A) \le \delta := \varepsilon/(2c)$,
\[
\E[\abs{X_i}\ind_A]
\le \E\bigl[\abs{X_i}\ind_{A \cap \{\abs{X_i}\le c\}}\bigr]
+ \E\bigl[\abs{X_i}\ind_{\{\abs{X_i}>c\}}\bigr]
\le c\,\Prob(A) + \frac{\varepsilon}{2} \le \varepsilon.
\]
\emph{(a), (b) $\Rightarrow$ UI.} Par l'inégalité de Markov,
$\Prob(\abs{X_i} > c) \le \E[\abs{X_i}]/c \le M/c$ avec
$M = \sup_i \E[\abs{X_i}]$, uniformément en $i$. Pour $\varepsilon$ donné,
prendre $\delta$ fourni par (b) puis $c > M/\delta$ : l'événement
$A_i = \{\abs{X_i} > c\}$ vérifie $\Prob(A_i) \le \delta$, donc
$\E[\abs{X_i}\ind_{A_i}] \le \varepsilon$ pour tout $i$.
\end{proof}

%-----------------------------------------------------------------------
\subsection{Critères pratiques}
%-----------------------------------------------------------------------

\begin{proposition}[Critères suffisants]\label{prop:criteres}
Chacune des conditions suivantes implique l'uniforme intégrabilité de
$(X_i)_{i\in I}$ :
\begin{enumerate}[label=(\roman*)]
\item \emph{bornitude dans $L^p$, $p > 1$} :
      $\sup_i \norm{X_i}_p < \infty$ ;
\item \emph{domination} : il existe $Y \in L^1$ telle que $\abs{X_i} \le Y$
      p.s.\ pour tout $i$.
\end{enumerate}
\end{proposition}

\begin{proof}
(i) Sur $\{\abs{X_i} > c\}$ on a $1 \le (\abs{X_i}/c)^{p-1}$, donc
\[
\E\bigl[\abs{X_i}\ind_{\{\abs{X_i}>c\}}\bigr]
\le \E\Bigl[\abs{X_i}\cdot \frac{\abs{X_i}^{p-1}}{c^{p-1}}\Bigr]
= \frac{\E[\abs{X_i}^p]}{c^{p-1}}
\le \frac{\sup_i \norm{X_i}_p^p}{c^{p-1}}
\xrightarrow[c\to\infty]{} 0.
\]
(ii) $\E[\abs{X_i}\ind_{\{\abs{X_i}>c\}}] \le \E[Y \ind_{\{Y>c\}}] \to 0$ par
convergence dominée, la borne étant indépendante de $i$.
\end{proof}

Le critère suivant est \emph{exact} : il caractérise l'UI.

\begin{theorem}[La Vallée Poussin]\label{thm:LVP}
$(X_i)_{i\in I}$ est UI si et seulement s'il existe une fonction
$\varphi : \R_+ \to \R_+$ croissante avec
$\varphi(x)/x \to \infty$ quand $x \to \infty$, telle que
\[
\sup_{i\in I} \E\bigl[\varphi(\abs{X_i})\bigr] < \infty.
\]
On peut de plus choisir $\varphi$ convexe.
\end{theorem}

\begin{proof}
\emph{Suffisance.} Posons $M = \sup_i \E[\varphi(\abs{X_i})]$ et
$\psi(c) = \inf_{x > c} \varphi(x)/x \to \infty$. Sur $\{\abs{X_i} > c\}$,
$\abs{X_i} \le \varphi(\abs{X_i})/\psi(c)$, donc
\[
\E\bigl[\abs{X_i}\ind_{\{\abs{X_i}>c\}}\bigr]
\le \frac{\E[\varphi(\abs{X_i})]}{\psi(c)} \le \frac{M}{\psi(c)}
\xrightarrow[c\to\infty]{} 0,
\]
uniformément en $i$.

\emph{Nécessité (esquisse).} Par UI, on peut choisir une suite
$c_k \uparrow \infty$ telle que
$\sup_i \E[\abs{X_i}\ind_{\{\abs{X_i}>c_k\}}] \le 2^{-k}$. On pose
$\varphi(x) = \sum_{k\ge1} (x - c_k)^+$, fonction convexe croissante ; comme le
nombre de termes actifs en $x$ tend vers l'infini, $\varphi(x)/x \to \infty$.
Enfin, en utilisant $(x-c)^+ \le x\,\ind_{\{x > c\}}$,
\[
\E\bigl[\varphi(\abs{X_i})\bigr]
\le \sum_{k\ge1} \E\bigl[\abs{X_i}\ind_{\{\abs{X_i}>c_k\}}\bigr]
\le \sum_{k\ge1} 2^{-k} = 1,
\]
uniformément en $i$.
\end{proof}

\begin{remark}
Le cas $\varphi(x) = x^p$ redonne le critère $L^p$ de la
proposition~\ref{prop:criteres}. Mais des croissances beaucoup plus lentes
suffisent : $\varphi(x) = x \log^+ x$ donne la classe $L\log L$, frontière
naturelle dans les inégalités maximales de Doob (l'inégalité
$\norm{\sup_n \abs{X_n}}_1 \lesssim 1 + \sup_n \E[\abs{X_n}\log^+\abs{X_n}]$
échoue dans $L^1$ pur) et en théorie ergodique.
\end{remark}

%-----------------------------------------------------------------------
\subsection{Le théorème de Vitali}
%-----------------------------------------------------------------------

L'UI est exactement le chaînon manquant entre convergence en probabilité et
convergence $L^1$.

\begin{theorem}[Vitali]\label{thm:vitali}
Soit $(X_n)$ une suite de variables intégrables convergeant en probabilité vers
$X$. Alors les assertions suivantes sont équivalentes :
\begin{enumerate}[label=(\roman*)]
\item $(X_n)$ est uniformément intégrable ;
\item $X \in L^1$ et $X_n \to X$ dans $L^1$.
\end{enumerate}
\end{theorem}

\begin{proof}
\emph{(i) $\Rightarrow$ (ii).} Une sous-suite $(X_{n_k})$ converge p.s.\ vers
$X$ ; par Fatou et la proposition~\ref{prop:UI-L1},
$\E[\abs{X}] \le \liminf_k \E[\abs{X_{n_k}}] < \infty$, donc $X \in L^1$.
Introduisons la troncature $\varphi_c(x) = (x \wedge c) \vee (-c)$, qui vérifie
$\abs{x - \varphi_c(x)} = (\abs{x} - c)^+ \le \abs{x}\ind_{\{\abs{x}>c\}}$.
Alors
\[
\E\abs{X_n - X}
\le \underbrace{\E\abs{X_n - \varphi_c(X_n)}}_{\le\, \sup_m \E[\abs{X_m}\ind_{\{\abs{X_m}>c\}}]}
+ \E\abs{\varphi_c(X_n) - \varphi_c(X)}
+ \underbrace{\E\abs{\varphi_c(X) - X}}_{\le\, \E[\abs{X}\ind_{\{\abs{X}>c\}}]}.
\]
Soit $\varepsilon > 0$. Par UI (premier terme) et intégrabilité de $X$
(troisième terme), on choisit $c$ rendant ces deux termes $\le \varepsilon$
uniformément en $n$. Le terme central tend vers $0$ : $\varphi_c$ étant
$1$-lipschitzienne, $\varphi_c(X_n) \to \varphi_c(X)$ en probabilité, et la
suite est bornée par $c$, donc la convergence a lieu dans $L^1$ (convergence
dominée le long de sous-suites p.s., puis argument de sous-sous-suites).
D'où $\limsup_n \E\abs{X_n - X} \le 2\varepsilon$.

\emph{(ii) $\Rightarrow$ (i).} On utilise la caractérisation de la
proposition~\ref{prop:UI-eps-delta}. La bornitude $L^1$ est claire. Pour
l'absolue continuité uniforme : soit $\varepsilon > 0$ et $N$ tel que
$\E\abs{X_n - X} \le \varepsilon/2$ pour $n > N$. La famille \emph{finie}
$\{X_1, \dots, X_N, X\}$ est UI (chaque variable intégrable l'est, et un
maximum fini de seuils convient), donc il existe $\delta$ tel que
$\Prob(A) \le \delta$ implique $\E[\abs{X_k}\ind_A] \le \varepsilon/2$ pour
$k \le N$ et $\E[\abs{X}\ind_A] \le \varepsilon/2$. Pour $n > N$,
$\E[\abs{X_n}\ind_A] \le \E\abs{X_n - X} + \E[\abs{X}\ind_A] \le \varepsilon$.
\end{proof}

\begin{remark}
Vitali généralise strictement la convergence dominée : la domination par un
$Y \in L^1$ fixe implique l'UI (proposition~\ref{prop:criteres}(ii)), mais l'UI
est plus faible, et surtout elle est \emph{nécessaire et suffisante}.
\end{remark}

%-----------------------------------------------------------------------
\subsection{Un fait structurel : les espérances conditionnelles d'une
variable fixée}
%-----------------------------------------------------------------------

\begin{theorem}\label{thm:cond-UI}
Soit $Z \in L^1$. La famille
\[
\bigl\{\E[Z \mid \G] \;:\; \G \text{ sous-tribu de } \F\bigr\}
\]
est uniformément intégrable.
\end{theorem}

\begin{proof}
Notons $X_\G = \E[Z \mid \G]$ et $c > 0$. Par l'inégalité de Jensen
conditionnelle, $\abs{X_\G} \le \E[\abs{Z} \mid \G]$, donc, l'événement
$\{\abs{X_\G} > c\}$ étant $\G$-mesurable,
\[
\E\bigl[\abs{X_\G}\,\ind_{\{\abs{X_\G}>c\}}\bigr]
\le \E\bigl[\E[\abs{Z}\mid\G]\,\ind_{\{\abs{X_\G}>c\}}\bigr]
= \E\bigl[\abs{Z}\,\ind_{\{\abs{X_\G}>c\}}\bigr],
\]
la dernière égalité par définition de l'espérance conditionnelle. Par Markov et
Jensen,
\[
\Prob\bigl(\abs{X_\G} > c\bigr) \le \frac{\E[\abs{X_\G}]}{c}
\le \frac{\E[\abs{Z}]}{c},
\]
majoration \emph{indépendante de $\G$}. Enfin, la mesure finie
$A \mapsto \E[\abs{Z}\ind_A]$ est absolument continue par rapport à $\Prob$ :
pour tout $\varepsilon$, il existe $\delta$ tel que $\Prob(A) \le \delta$
implique $\E[\abs{Z}\ind_A] \le \varepsilon$ (tronquer $\abs{Z}$ à un niveau
$c_0$ tel que $\E[\abs{Z}\ind_{\{\abs{Z}>c_0\}}] \le \varepsilon/2$, puis
$\delta = \varepsilon/(2c_0)$). En combinant : pour $c \ge \E[\abs{Z}]/\delta$,
tous les événements $\{\abs{X_\G}>c\}$ ont probabilité $\le \delta$, donc
$\sup_\G \E[\abs{X_\G}\ind_{\{\abs{X_\G}>c\}}] \le \varepsilon$.
\end{proof}

Ce résultat est remarquable : \emph{le conditionnement ne peut pas créer de
fuite de masse}, quelle que soit la tribu. Il fournit immédiatement la moitié
facile de la caractérisation des martingales fermées.

%=======================================================================
\section{Martingales fermées et théorème d'arrêt}
%=======================================================================

%-----------------------------------------------------------------------
\subsection{La caractérisation}
%-----------------------------------------------------------------------

\begin{theorem}[Martingales fermées]\label{thm:fermee}
Soit $(X_n)_{n\ge0}$ une martingale. Les assertions suivantes sont
équivalentes :
\begin{enumerate}[label=(\roman*)]
\item la famille $(X_n)_{n \ge 0}$ est uniformément intégrable ;
\item $(X_n)$ converge dans $L^1$ ;
\item il existe $Z \in L^1$ telle que $X_n = \E[Z \mid \F_n]$ pour tout $n$
      (la martingale est \emph{fermée}).
\end{enumerate}
Dans ce cas, $X_n \to X_\infty$ p.s.\ et dans $L^1$, et
$X_n = \E[X_\infty \mid \F_n]$ pour tout $n$ ; on peut prendre
$X_\infty = \E[Z \mid \F_\infty]$ où $\F_\infty = \sigma\bigl(\bigcup_n
\F_n\bigr)$.
\end{theorem}

\begin{proof}
\emph{(i) $\Rightarrow$ (ii).} L'UI donne la bornitude $L^1$
(proposition~\ref{prop:UI-L1}), donc la convergence p.s.\ vers $X_\infty$ par le
théorème~\ref{thm:doob} ; la convergence p.s.\ implique la convergence en
probabilité, et Vitali (théorème~\ref{thm:vitali}) conclut à la convergence
$L^1$.

\emph{(ii) $\Rightarrow$ (iii).} Notons $X_\infty$ la limite $L^1$. Fixons $n$
et faisons tendre $m \to \infty$ dans l'identité de martingale
$X_n = \E[X_m \mid \F_n]$ ($m \ge n$) : l'espérance conditionnelle est une
contraction de $L^1$
($\E\abs{\E[X_m \mid \F_n] - \E[X_\infty \mid \F_n]} \le \E\abs{X_m -
X_\infty}$), donc $X_n = \E[X_\infty \mid \F_n]$, et (iii) tient avec
$Z = X_\infty$.

\emph{(iii) $\Rightarrow$ (i).} C'est le théorème~\ref{thm:cond-UI} appliqué à
la sous-famille $\{\E[Z\mid\F_n]\}_{n\ge0}$.

\emph{Identification de la limite.} Supposons (iii). La martingale est UI donc
converge p.s.\ et $L^1$ vers une limite $X_\infty$, qui est
$\F_\infty$-mesurable. Montrons $X_\infty = \E[Z \mid \F_\infty]$. Pour
$A \in \F_n$ fixé et $m \ge n$,
$\E[X_m \ind_A] = \E\bigl[\E[Z\mid\F_m]\ind_A\bigr] = \E[Z\ind_A]$ ; la
convergence $L^1$ donne à la limite $\E[X_\infty \ind_A] = \E[Z \ind_A]$. Les
deux mesures finies signées $A \mapsto \E[X_\infty\ind_A]$ et
$A \mapsto \E[Z\ind_A]$ coïncident donc sur le $\pi$-système
$\bigcup_n \F_n$, stable par intersection finie et engendrant $\F_\infty$ ; par
le lemme de classe monotone, elles coïncident sur $\F_\infty$, ce qui
caractérise $X_\infty = \E[Z\mid\F_\infty]$.
\end{proof}

\begin{corollary}[Théorème de convergence de Lévy]\label{cor:levy}
Pour toute $Z \in L^1$,
\[
\E[Z \mid \F_n] \xrightarrow[n\to\infty]{\ps \text{ et } L^1}
\E[Z \mid \F_\infty].
\]
\end{corollary}

\begin{remark}[Deux mondes]\label{rem:deux-mondes}
Le théorème~\ref{thm:fermee} trace une frontière structurelle au sein des
martingales convergentes.
\begin{itemize}
\item \emph{Martingales fermées} : $X_n = \E[X_\infty \mid \F_n]$. La
martingale est la révélation progressive d'une variable terminale ; toute
l'information vit \og au bout \fg{} et le processus se reconstruit depuis sa
limite. C'est le cas de l'urne de Pólya ($M_n = \E[M_\infty \mid \F_n]$ avec
$M_\infty$ de loi Beta) et du processus de densité dans le cas équivalent de la
dichotomie de Kakutani.
\item \emph{Martingales non fermées} : la limite p.s.\ existe (si bornitude
$L^1$) mais a \og oublié \fg{} de la masse. Pour une martingale positive, le
lemme de Fatou donne seulement $\E[X_\infty] \le \liminf \E[X_n]$, et
l'inégalité peut être stricte. Le processus n'est alors \emph{pas} déterminé
par sa limite : dans le cas dégénéré de Kakutani, $\E[X_\infty \mid \F_n] = 0
\ne X_n$.
\end{itemize}
La dichotomie de Kakutani (section~\ref{sec:kakutani}) est précisément un
critère \emph{exact} de fermeture pour la classe des martingales produit.
\end{remark}

%-----------------------------------------------------------------------
\subsection{Le théorème d'arrêt sous uniforme intégrabilité}
%-----------------------------------------------------------------------

Rappelons que pour un temps d'arrêt $T$, la tribu du passé jusqu'à $T$ est
$\F_T = \{A \in \F : A \cap \{T \le n\} \in \F_n \ \forall n\}$.

\begin{theorem}[Arrêt optionnel, version UI]\label{thm:arret}
Soit $(X_n)$ une martingale uniformément intégrable, de limite $X_\infty$, et
soient $S \le T$ deux temps d'arrêt quelconques (éventuellement infinis, en
posant $X_T = X_\infty$ sur $\{T = \infty\}$). Alors $X_S, X_T \in L^1$ et
\[
\E[X_T \mid \F_S] = X_S \qquad \ps ;
\qquad \text{en particulier } \E[X_T] = \E[X_S] = \E[X_0].
\]
\end{theorem}

\begin{proof}
\emph{Étape 1 : $X_T = \E[X_\infty \mid \F_T]$.} Soit $A \in \F_T$. On
décompose $A = \bigsqcup_{n \in \N \cup \{\infty\}} A \cap \{T = n\}$ avec
$A \cap \{T=n\} \in \F_n$. Par fermeture
(théorème~\ref{thm:fermee}), $X_n = \E[X_\infty \mid \F_n]$, donc
\[
\E\bigl[X_T \ind_{A \cap \{T=n\}}\bigr]
= \E\bigl[X_n \ind_{A \cap \{T=n\}}\bigr]
= \E\bigl[X_\infty \ind_{A \cap \{T=n\}}\bigr].
\]
La sommation sur $n$ est licite : les termes
$\E[\abs{X_\infty}\ind_{A\cap\{T=n\}}]$ sont sommables (de somme
$\le \E\abs{X_\infty}$), et pour le membre de gauche on utilise l'UI de la
famille $\{\E[X_\infty\mid\G]\}$ (théorème~\ref{thm:cond-UI}) qui rend la série
absolument convergente. D'où $\E[X_T\ind_A] = \E[X_\infty\ind_A]$ pour tout
$A \in \F_T$, ce qui est exactement $X_T = \E[X_\infty\mid\F_T]$ ; en
particulier $X_T \in L^1$.

\emph{Étape 2 : tour d'espérances conditionnelles.} Comme $S \le T$, on a
$\F_S \subseteq \F_T$, et par emboîtement
\[
\E[X_T \mid \F_S] = \E\bigl[\E[X_\infty \mid \F_T] \,\big|\, \F_S\bigr]
= \E[X_\infty \mid \F_S] = X_S. \qedhere
\]
\end{proof}

\begin{example}[La stratégie de doublement]\label{ex:doubling}
Sans UI, le théorème échoue même pour $T < \infty$ p.s. Soient $(\varepsilon_k)$
i.i.d.\ uniformes sur $\{-1,+1\}$ et la stratégie de mise doublée : on mise
$2^{k-1}$ à la $k$-ième partie tant qu'on a perdu, on s'arrête au premier gain,
$T = \inf\{k : \varepsilon_k = +1\}$ (fini p.s., de loi géométrique). Le
capital $X_n = \sum_{k=1}^{n \wedge T} 2^{k-1}\varepsilon_k$ est une
transformée de martingale, donc une martingale (lemme~\ref{lem:transform}),
avec $\E[X_n] = 0$. Sur $\{T \le n\}$, les pertes $1 + 2 + \cdots + 2^{T-2} =
2^{T-1} - 1$ sont exactement compensées et dépassées : $X_n = 1$. Sur
$\{T > n\}$ (probabilité $2^{-n}$), $X_n = -(2^n - 1)$. Ainsi
\[
X_n \xrightarrow{\ps} X_T = 1, \qquad \text{mais} \qquad
\E[X_T] = 1 \ne 0 = \E[X_0].
\]
La martingale est pourtant bornée dans $L^1$
($\E\abs{X_n} = (1 - 2^{-n}) + 2^{-n}(2^n - 1) \le 2$) : c'est l'UI qui fait
défaut, la masse négative $-(2^n-1) \cdot 2^{-n} \approx -1$ étant portée par
des valeurs divergentes --- de nouveau le mécanisme de fuite à l'infini.
\end{example}

%=======================================================================
\section{La dichotomie de Kakutani}\label{sec:kakutani}
%=======================================================================

%-----------------------------------------------------------------------
\subsection{Formulation martingale}
%-----------------------------------------------------------------------

Le critère de Kakutani s'exprime par un produit infini. Commençons par lever
toute ambiguïté sur le sens de ce produit : pour des facteurs dans $(0,1]$, la
limite des produits partiels existe \emph{toujours}, et la seule question est
de savoir si elle est nulle ou strictement positive.

\begin{lemma}[{Produits infinis de facteurs dans $(0,1]$}]\label{lem:prod-inf}
Soit $(a_k)_{k\ge1}$ une suite de réels avec $0 < a_k \le 1$ pour tout $k$, et
$\pi_n = \prod_{k=1}^{n} a_k$.
\begin{enumerate}[label=(\roman*)]
\item La suite $(\pi_n)$ converge toujours : décroissante et minorée par $0$,
      elle admet une limite $\pi_\infty \in [0,1]$. C'est cette limite que
      l'on note $\prod_{k\ge1} a_k$.
\item On a l'équivalence
\[
\pi_\infty > 0
\iff \sum_{k\ge1} (-\log a_k) < \infty
\iff \sum_{k\ge1} (1 - a_k) < \infty.
\]
\end{enumerate}
\end{lemma}

\begin{proof}
(i) Chaque facteur étant dans $(0,1]$, la suite $(\pi_n)$ est décroissante,
strictement positive, minorée par $0$ : elle converge vers
$\pi_\infty \in [0,1]$. Aucune autre hypothèse n'est requise --- contrairement
au cas de facteurs de signe quelconque, où l'existence de la limite est déjà
un problème, la monotonie fait ici tout le travail. (Dans le vocabulaire
classique des produits infinis, le cas $\pi_\infty = 0$ est appelé
\emph{divergence vers $0$} ; nous n'utilisons pas cette convention et parlons
simplement de la valeur de la limite.)

(ii) \emph{Première équivalence.} Comme $a_k > 0$, on peut écrire
$\log \pi_n = \sum_{k\le n} \log a_k$, somme décroissante de termes $\le 0$ ;
elle est minorée si et seulement si $\sum_k(-\log a_k) < \infty$, et
$\pi_\infty = \exp\bigl(-\sum_k(-\log a_k)\bigr) > 0$ exactement dans ce cas.

\emph{Seconde équivalence.} On utilise deux inégalités élémentaires, toutes
deux issues de $\log x \le x - 1$ :
\[
1 - x \le -\log x \quad (x \in (0,1]),
\qquad\qquad
-\log x = \log\frac1x \le \frac1x - 1 = \frac{1-x}{x} \le 2(1-x)
\quad (x \in [\tfrac12, 1]).
\]
Si $\sum_k(-\log a_k) < \infty$, la première inégalité donne
$\sum_k(1-a_k) < \infty$. Réciproquement, si $\sum_k (1-a_k) < \infty$, alors
$1 - a_k \to 0$, donc $a_k \ge \tfrac12$ à partir d'un rang $k_0$ ; la seconde
inégalité donne
$\sum_{k \ge k_0}(-\log a_k) \le 2\sum_{k\ge k_0}(1-a_k) < \infty$, et les
termes d'indice $< k_0$, en nombre fini, ne changent rien.

Enfin, si $a_k \not\to 1$ (cas où aucune des deux séries ne converge), la
conclusion $\pi_\infty = 0$ se voit aussi directement : il existe $\delta > 0$
et une infinité d'indices $k$ avec $a_k \le 1 - \delta$, d'où
$\pi_n \le (1-\delta)^{N_n}$ avec $N_n \to \infty$.
\end{proof}

\begin{theorem}[Kakutani, 1948 --- version martingale]\label{thm:kakutani-mart}
Soit $(\xi_k)_{k \ge 1}$ une suite de variables aléatoires indépendantes,
positives, d'espérance $\E[\xi_k] = 1$, et soit
\[
X_n = \prod_{k=1}^{n} \xi_k, \qquad X_0 = 1,
\]
la martingale produit associée. On pose
$a_k = \E\bigl[\sqrt{\xi_k}\bigr] \in (0,1]$ ; d'après le
lemme~\ref{lem:prod-inf}, le produit $\prod_{k\ge1} a_k$, entendu comme limite
des produits partiels, existe toujours dans $[0,1]$. Alors exactement l'une des
deux situations suivantes se produit :
\begin{enumerate}[label=(\Alph*)]
\item $\displaystyle\prod_{k\ge 1} a_k > 0$
      \emph{(de manière équivalente : $\sum_k (1-a_k) < \infty$)} :
      la martingale $(X_n)$ est uniformément intégrable, $X_n \to X_\infty$
      p.s.\ et dans $L^1$, $\E[X_\infty] = 1$, et $\Prob(X_\infty > 0) = 1$ ;
\item $\displaystyle\prod_{k\ge 1} a_k = 0$ :
      \quad $X_\infty = 0$ presque sûrement.
\end{enumerate}
\end{theorem}

\begin{remark}
(a) Vérifions que $a_k \in (0,1]$, ce qui rend le lemme~\ref{lem:prod-inf}
applicable et donne un sens inconditionnel au produit $\prod_k a_k$.
La majoration $a_k \le 1$ résulte de Cauchy--Schwarz :
$\E[\sqrt{\xi_k}] \le \sqrt{\E[\xi_k]} = 1$, avec égalité si et seulement si
$\xi_k = 1$ p.s. La minoration $a_k > 0$ vient de ce que $\E[\xi_k] = 1$
impose $\Prob(\xi_k > 0) > 0$, donc $\E[\sqrt{\xi_k}] > 0$.

(b) L'équivalence entre $\prod_k a_k > 0$ et $\sum_k (1 - a_k) < \infty$ est
exactement le point (ii) du lemme~\ref{lem:prod-inf}.

(c) Le produit $\prod_k a_k$ mesure l'accumulation des perturbations : chaque
facteur non trivial coûte strictement, et la question est de savoir si le coût
total est fini.
\end{remark}

\begin{proof}
Dans les deux cas, $(X_n)$ est une martingale positive, donc converge p.s.\ vers
une limite $X_\infty \ge 0$ intégrable (corollaire~\ref{cor:cas-part}(iii)).

\medskip
\noindent\emph{Cas (A).} Posons $\pi_n = \prod_{k=1}^n a_k$ et
\[
Y_n = \prod_{k=1}^{n} \frac{\sqrt{\xi_k}}{a_k}.
\]
Les facteurs $\sqrt{\xi_k}/a_k$ sont indépendants, positifs, d'espérance $1$,
donc $(Y_n)$ est une martingale positive de moyenne $1$. Son moment d'ordre
deux se calcule par indépendance :
\[
\E\bigl[Y_n^2\bigr]
= \prod_{k=1}^{n} \frac{\E[\xi_k]}{a_k^2}
= \prod_{k=1}^{n} a_k^{-2}
= \pi_n^{-2}
\;\le\; \Bigl(\prod_{k\ge 1} a_k\Bigr)^{-2} =: K < \infty.
\]
La martingale $(Y_n)$ est bornée dans $L^2$, donc UI
(proposition~\ref{prop:criteres}(i)) ; elle converge p.s., et la convergence a
lieu dans $L^2$ : en effet, pour $m < n$, l'orthogonalité des accroissements de
martingale dans $L^2$ donne
$\E[(Y_n - Y_m)^2] = \E[Y_n^2] - \E[Y_m^2] = \pi_n^{-2} - \pi_m^{-2} \to 0$,
la suite $(\pi_n^{-2})$ étant croissante convergente : $(Y_n)$ est de Cauchy
dans $L^2$.

Comme $X_n = Y_n^2\, \pi_n^2$, il reste à voir que $Y_n^2 \to Y_\infty^2$ dans
$L^1$. Par Cauchy--Schwarz,
\[
\E\bigl[\abs{Y_n^2 - Y_\infty^2}\bigr]
= \E\bigl[\abs{Y_n - Y_\infty}\,(Y_n + Y_\infty)\bigr]
\le \norm{Y_n - Y_\infty}_2 \,\norm{Y_n + Y_\infty}_2
\le 2\sqrt{K}\,\norm{Y_n - Y_\infty}_2 \to 0.
\]
Avec $\pi_n^2 \to \pi_\infty^2 > 0$, on conclut que $X_n \to X_\infty =
Y_\infty^2 \pi_\infty^2$ dans $L^1$. Le théorème~\ref{thm:fermee} donne l'UI et
$\E[X_\infty] = \lim_n \E[X_n] = 1$.

Montrons enfin $\Prob(X_\infty > 0) = 1$ et non seulement $> 0$. Supposons de
plus $\xi_k > 0$ p.s.\ pour tout $k$ --- hypothèse naturelle, puisque dans
l'application aux mesures produit les $\xi_k$ sont des densités de mesures
équivalentes ; elle est d'ailleurs nécessaire, car $\Prob(\xi_{k_0} = 0) > 0$
pour un seul $k_0$ entraîne $\Prob(X_\infty = 0) > 0$. Sous cette hypothèse,
modifier un nombre fini de facteurs $\xi_1, \dots, \xi_m$ multiplie $X_\infty$
par une constante aléatoire strictement positive, donc l'événement
$\{X_\infty > 0\}$ appartient à la tribu asymptotique
$\bigcap_m \sigma(\xi_k,\, k > m)$. Par la loi du zéro--un de Kolmogorov,
$\Prob(X_\infty > 0) \in \{0, 1\}$, et $\E[X_\infty] = 1$ exclut la valeur
$0$ : ainsi $\Prob(X_\infty > 0) = 1$.

\medskip
\noindent\emph{Cas (B).} Par indépendance,
\[
\E\bigl[\sqrt{X_n}\bigr] = \prod_{k=1}^{n} \E[\sqrt{\xi_k}]
= \pi_n \xrightarrow[n\to\infty]{} 0.
\]
Comme $\sqrt{X_n} \to \sqrt{X_\infty}$ p.s.\ (continuité de la racine), le
lemme de Fatou donne
\[
\E\bigl[\sqrt{X_\infty}\bigr] \le \liminf_n \E\bigl[\sqrt{X_n}\bigr] = 0,
\]
donc $\sqrt{X_\infty} = 0$ p.s., soit $X_\infty = 0$ p.s.
\end{proof}

\begin{remark}[Pourquoi la racine carrée ?]
Le choix de l'exposant $1/2$ n'est pas anecdotique. Pour $\theta \in (0,1)$, la
fonction $x \mapsto x^\theta$ est concave, donc $\E[\xi_k^\theta] \le 1$ et le
même argument de Fatou fonctionne dans le cas dégénéré. Mais c'est $\theta =
1/2$ qui rend le calcul du cas (A) exact : $\E[(\sqrt{\xi_k}/a_k)^2] =
\E[\xi_k]/a_k^2 = a_k^{-2}$ utilise précisément le fait que le carré de la
racine redonne $\xi_k$, d'espérance connue. L'exposant $1/2$ est aussi celui
qui fait apparaître l'affinité de Hellinger, quantité symétrique en $(P,Q)$ ---
voir ci-dessous.
\end{remark}

%-----------------------------------------------------------------------
\subsection{Formulation en théorie de la mesure}
%-----------------------------------------------------------------------

\begin{theorem}[Kakutani --- mesures produit]\label{thm:kakutani-mes}
Soient $(E_k, \mathcal{E}_k)_{k\ge1}$ des espaces mesurables et, sur chacun,
deux probabilités équivalentes $P_k \sim Q_k$. On définit sur l'espace produit
$E = \prod_k E_k$, muni de la tribu produit, les mesures
\[
P = \bigotimes_{k\ge1} P_k, \qquad Q = \bigotimes_{k\ge1} Q_k,
\]
et l'\emph{affinité de Hellinger} de chaque facteur
\[
\rho_k = \int_{E_k} \sqrt{\frac{dQ_k}{dP_k}}\; dP_k
= \int_{E_k} \sqrt{\frac{dQ_k}{d\mu_k}\cdot\frac{dP_k}{d\mu_k}}\; d\mu_k
\;\in\; (0,1],
\]
(la seconde écriture, pour toute mesure dominante $\mu_k$, montre la symétrie
en $(P_k, Q_k)$). Alors :
\begin{enumerate}[label=(\Alph*)]
\item si $\prod_k \rho_k > 0$, alors $Q \sim P$ (mesures équivalentes), avec
      $\frac{dQ}{dP} = X_\infty$ où $X_\infty$ est la limite du processus de
      densité ;
\item si $\prod_k \rho_k = 0$, alors $Q \perp P$ (mesures étrangères).
\end{enumerate}
\end{theorem}

\begin{proof}
Travaillons sous $P$, avec les coordonnées canoniques
$\omega = (\omega_1, \omega_2, \dots)$, la filtration
$\F_n = \sigma(\omega_1, \dots, \omega_n)$ et
\[
\xi_k(\omega) = \frac{dQ_k}{dP_k}(\omega_k).
\]
Sous $P$, les $\xi_k$ sont indépendantes, strictement positives
($P_k \sim Q_k$), d'espérance $\E_P[\xi_k] = \int \frac{dQ_k}{dP_k}\,dP_k = 1$,
et $\E_P[\sqrt{\xi_k}] = \rho_k$. La martingale produit
$X_n = \prod_{k\le n}\xi_k$ est le \emph{processus de densité} : pour tout
$A \in \F_n$ (dépendant des $n$ premières coordonnées), le théorème de Fubini
donne
\[
Q(A) = \E_P\bigl[X_n \ind_A\bigr],
\qquad \text{i.e.} \qquad
X_n = \frac{dQ}{dP}\bigg|_{\F_n}.
\]

\emph{Cas (A).} Par le théorème~\ref{thm:kakutani-mart}, $X_n \to X_\infty$
dans $L^1(P)$ avec $X_\infty > 0$ $P$-p.s. Pour $A \in \F_n$ fixé, la
convergence $L^1$ permet de passer à la limite :
\[
Q(A) = \lim_{m\to\infty}\E_P[X_m \ind_A] = \E_P[X_\infty \ind_A].
\]
Les deux mesures finies $Q$ et $X_\infty\,dP$ coïncident sur le $\pi$-système
$\bigcup_n \F_n$ qui engendre la tribu produit ; par le lemme de classe
monotone, $dQ = X_\infty\, dP$ sur toute la tribu produit. Comme $X_\infty > 0$
$P$-p.s., les mesures sont équivalentes : $Q \sim P$.

\emph{Cas (B).} Sous $P$, $X_\infty = 0$ p.s. Par symétrie de $\rho_k$,
l'argument s'applique aussi sous $Q$ au processus de densité inverse
$\widetilde{X}_n = \prod_{k \le n} \frac{dP_k}{dQ_k}(\omega_k) = 1/X_n$ : on a
$\E_Q[\sqrt{\widetilde\xi_k}] = \int\sqrt{dP_k/dQ_k}\,dQ_k = \rho_k$, donc
$\widetilde{X}_n \to 0$ $Q$-p.s., c'est-à-dire $X_n \to +\infty$ $Q$-p.s.
L'événement asymptotique
\[
D = \Bigl\{\omega : \lim_n X_n(\omega) = 0\Bigr\}
\]
vérifie donc $P(D) = 1$ et $Q(D) = 0$ : les deux mesures se concentrent sur des
événements disjoints ($D$ et $D^c$), soit $Q \perp P$.
\end{proof}

\begin{remark}[Tout ou rien]
Le contenu frappant du théorème est l'absence de décomposition de Lebesgue non
triviale : bien que chaque facteur soit équivalent, une infinité de
perturbations individuellement bénignes produit soit une équivalence globale,
soit une singularité totale --- jamais un mélange. La distance de Hellinger
$H^2(P_k, Q_k) = \int(\sqrt{dP_k} - \sqrt{dQ_k})^2 = 2(1 - \rho_k)$ reformule
le critère : équivalence si et seulement si $\sum_k H^2(P_k, Q_k) < \infty$.
\end{remark}

%=======================================================================
\section{Le cas gaussien et le théorème de Cameron--Martin}
%=======================================================================

\begin{example}[Martingale exponentielle gaussienne]\label{ex:gauss}
Soient $(Z_k)$ i.i.d.\ $\mathcal{N}(0,1)$, $(\sigma_k)$ une suite réelle, et
\[
\xi_k = \exp\Bigl(\sigma_k Z_k - \tfrac{\sigma_k^2}{2}\Bigr),
\qquad
X_n = \prod_{k=1}^n \xi_k
= \exp\Bigl(\sum_{k=1}^n \sigma_k Z_k - \tfrac12 \sum_{k=1}^n \sigma_k^2\Bigr).
\]
\emph{Calculs exacts.} La transformée de Laplace gaussienne
$\E[e^{\lambda Z}] = e^{\lambda^2/2}$ donne d'une part
$\E[\xi_k] = e^{\sigma_k^2/2 - \sigma_k^2/2} = 1$ (martingale), d'autre part
\[
a_k = \E\bigl[\sqrt{\xi_k}\bigr]
= \E\Bigl[e^{\frac{\sigma_k}{2} Z_k}\Bigr]\, e^{-\frac{\sigma_k^2}{4}}
= e^{\frac{\sigma_k^2}{8}}\, e^{-\frac{\sigma_k^2}{4}}
= e^{-\sigma_k^2/8}.
\]
La dichotomie se lit alors sur une condition de sommabilité :
\[
\prod_k a_k = \exp\Bigl(-\tfrac18\textstyle\sum_k \sigma_k^2\Bigr) > 0
\iff
\boxed{\;\sum_{k\ge1} \sigma_k^2 < \infty.\;}
\]

\emph{Régime convergent.} Si $\sum_k \sigma_k^2 < \infty$, la série
$\sum_k \sigma_k Z_k$ converge p.s.\ et dans $L^2$ (série de variables
indépendantes centrées de variances sommables, théorème des deux séries de
Kolmogorov), donc
\[
X_\infty = \exp\Bigl(\sum_{k\ge1}\sigma_k Z_k
- \tfrac12\sum_{k\ge1}\sigma_k^2\Bigr) > 0 \quad \ps,
\qquad \E[X_\infty] = 1 :
\]
la martingale est fermée, de limite log-normale explicite.

\emph{Régime dégénéré.} Si $\sum_k \sigma_k^2 = \infty$, alors $X_\infty = 0$
p.s.\ bien que $\E[X_n] = 1$ pour tout $n$. À $\sigma_k = \sigma > 0$ constant,
la loi des grands nombres précise la vitesse :
\[
\frac1n \log X_n = \sigma\,\frac{1}{n}\sum_{k=1}^n Z_k - \frac{\sigma^2}{2}
\xrightarrow{\ps} -\frac{\sigma^2}{2} < 0,
\]
décroissance exponentielle p.s. L'espérance, elle, est portée par un événement
de grande déviation : $X_n \ge 1$ exige
$\bar{Z}_n := \frac1n\sum_{k\le n} Z_k \ge \sigma/2$, événement de probabilité
\[
\Prob\bigl(\bar{Z}_n \ge \sigma/2\bigr)
= \Prob\bigl(\mathcal{N}(0,1) \ge \tfrac{\sigma}{2}\sqrt{n}\bigr)
\sim \frac{e^{-n\sigma^2/8}}{\frac{\sigma}{2}\sqrt{2\pi n}},
\]
exponentiellement petit --- et le taux $\sigma^2/8$ est exactement
$-\log a_k$ : la vitesse de Hellinger \emph{est} le taux de grande déviation de
l'événement qui porte la masse. C'est ce mécanisme qui rend un estimateur de
Monte-Carlo de $\E[X_n]$ inutilisable à $n$ grand : la masse fuit le long
d'événements jamais échantillonnés.
\end{example}

\begin{corollary}[Cameron--Martin, cas discret]\label{cor:CM}
Sur $\R^{\N}$ muni de la tribu produit,
\[
\bigotimes_{k\ge1} \mathcal{N}(\theta_k, 1) \;\sim\; \bigotimes_{k\ge1}
\mathcal{N}(0,1)
\iff
(\theta_k)_k \in \ell^2,
\]
et il y a singularité totale ($\perp$) sinon. Dans le cas équivalent, la
densité est
\[
\frac{d\bigotimes_k \mathcal{N}(\theta_k,1)}{d\bigotimes_k \mathcal{N}(0,1)}(x)
= \exp\Bigl(\sum_{k\ge1}\theta_k x_k - \tfrac12\sum_{k\ge1}\theta_k^2\Bigr).
\]
\end{corollary}

\begin{proof}
La densité du facteur $k$ est
\[
\frac{d\mathcal{N}(\theta_k,1)}{d\mathcal{N}(0,1)}(x)
= \frac{e^{-(x-\theta_k)^2/2}}{e^{-x^2/2}}
= e^{\theta_k x - \theta_k^2/2},
\]
soit exactement le $\xi_k$ de l'exemple~\ref{ex:gauss} avec
$\sigma_k = \theta_k$ et $Z_k = x_k$ sous la mesure de référence. Le calcul
$\rho_k = e^{-\theta_k^2/8}$ et le théorème~\ref{thm:kakutani-mes} concluent :
équivalence si et seulement si $\sum_k \theta_k^2 < \infty$. La formule de la
densité est la limite $X_\infty$ du régime convergent.
\end{proof}

\begin{remark}
C'est la version discrète du théorème de Cameron--Martin pour la mesure de
Wiener : la translation du brownien par une fonction $h$,
$W \mapsto W + h$, produit une mesure équivalente si et seulement si $h$
appartient à l'espace de Cameron--Martin $H^1 = \{h : h' \in L^2\}$, avec
densité de Girsanov
$\exp\bigl(\int_0^\infty h'\,dW - \frac12\int_0^\infty (h')^2\,dt\bigr)$, et
une mesure étrangère sinon. En décomposant sur une base hilbertienne, on
retrouve exactement le corollaire~\ref{cor:CM} : même mécanisme, même
dichotomie, avec $\sum \theta_k^2$ remplacée par $\norm{h'}_{L^2}^2$.
\end{remark}

%=======================================================================
\section{Portée du résultat}
%=======================================================================

\paragraph{Statistique asymptotique.}
Soit à tester $H_0$ : les observations suivent $P$, contre $H_1$ : elles
suivent $Q$, à partir de la suite infinie d'observations indépendantes. La
dichotomie affirme que la discrimination parfaite est soit possible (cas
singulier : le test de région critique $\{X_\infty = 0\}^c$, i.e.\ rejeter
$H_0$ lorsque le rapport de vraisemblance ne tend pas vers $0$, ne commet
asymptotiquement aucune erreur de première ni de seconde espèce), soit
impossible même asymptotiquement (cas équivalent : aucun test consistant
n'existe, les deux suites d'hypothèses restent \emph{contiguës} au sens de
Le~Cam). L'affinité de Hellinger $\rho_k$ et la distance associée
$H^2(P_k, Q_k) = 2(1 - \rho_k)$ sont ainsi les quantités naturelles de la
théorie : la frontière $\sum_k H^2(P_k, Q_k) < \infty$ sépare exactement le
détectable de l'indétectable. Dans le cas gaussien du
corollaire~\ref{cor:CM}, cette frontière est le critère $\ell^2$ sur le signal
$(\theta_k)$ --- forme primitive des seuils de détection en tests de signaux
parcimonieux.

\paragraph{Finance mathématique.}
Le processus de densité d'un changement de mesure de Girsanov,
\[
X_t = \exp\Bigl(\int_0^t \theta_s\, dW_s - \tfrac12 \int_0^t \theta_s^2\,
ds\Bigr),
\]
est une martingale exponentielle du type de l'exemple~\ref{ex:gauss} (une
martingale locale positive, donc une surmartingale, dont il faut vérifier
qu'elle est une vraie martingale). Sur horizon fini $T$, la condition de
Novikov $\E\bigl[e^{\frac12\int_0^T \theta_s^2\, ds}\bigr] < \infty$ garantit
que $(X_t)_{t \le T}$ est une martingale UI fermée par $X_T$, et la mesure
risque-neutre $dQ = X_T\, dP$ existe. Sur horizon infini, c'est le mécanisme de
Kakutani qui gouverne : si la prime de risque cumulée
$\int_0^\infty \theta_s^2\, ds$ diverge, le changement de mesure dégénère
($X_\infty = 0$), et les mesures historique et risque-neutre deviennent
étrangères --- même mathématique que le cas discret, avec $\sum_k \sigma_k^2$
remplacée par l'intégrale. La stratégie de doublement de
l'exemple~\ref{ex:doubling} est l'autre visage financier du défaut d'UI :
c'est précisément pour l'exclure que la théorie de l'arbitrage impose des
contraintes d'admissibilité (richesse minorée) aux stratégies.

\paragraph{En résumé.}
Le théorème de Doob garantit l'existence de la limite ; l'uniforme
intégrabilité est la condition exacte sous laquelle la martingale est \emph{ce
qu'elle prétend être} --- l'espérance conditionnelle de sa variable terminale
--- et sous laquelle les identités formelles ($\E[X_T] = \E[X_0]$, passage à la
limite dans le conditionnement) deviennent licites ; la dichotomie de Kakutani
répond à la question que Doob laisse ouverte --- \emph{la limite
conserve-t-elle la masse ?} --- et y répond de la manière la plus tranchée
possible : tout ou rien, avec un critère exact, calculable facteur par facteur.

\begin{thebibliography}{9}
\bibitem{kakutani} S.~Kakutani, \emph{On equivalence of infinite product
measures}, Annals of Mathematics \textbf{49} (1948), 214--224.
\bibitem{williams} D.~Williams, \emph{Probability with Martingales}, Cambridge
University Press, 1991.
\bibitem{durrett} R.~Durrett, \emph{Probability: Theory and Examples},
5\textsuperscript{e} éd., Cambridge University Press, 2019.
\bibitem{lecam} L.~Le~Cam, \emph{Asymptotic Methods in Statistical Decision
Theory}, Springer, 1986.
\bibitem{revuzyor} D.~Revuz et M.~Yor, \emph{Continuous Martingales and
Brownian Motion}, 3\textsuperscript{e} éd., Springer, 1999.
\end{thebibliography}

\end{document}
